\Askhsh{27667_SOLUTION}{1}
\begin{ylh}{1}
    \item [Συναρτήσεις]
\end{ylh}

\Askhseis{\alpha)}
\begin{rlist}
    \item Η συνάρτηση $f$ είναι συνεχής στο $\mathbb{R}$ (ως άθροισμα πολυωνυμικής με εκθετική). Για κάθε $x \in \mathbb{R}$ έχουμε:
    \[ f'(x) = \left(e^x + \frac{x^2}{2} + 2023\right)' = e^x + x \]
    \[ f''(x) = (e^x + x)' = e^x + 1 > 0 \]
    άρα η συνάρτηση $f$ είναι κυρτή στο $\mathbb{R}$.
    \item Επειδή η συνάρτηση $f$ είναι κυρτή στο $\mathbb{R}$ έπεται ότι η συνάρτηση $f'$ είναι γνησίως αύξουσα στο $\mathbb{R}$. Ακόμη η $f'$ είναι και συνεχής ως παραγωγίσιμη.

    Οπότε $f'(\mathbb{R}) = \left(\lim_{x \to -\infty} f'(x), \lim_{x \to +\infty} f'(x)\right) = \mathbb{R}$, αφού είναι:
    \begin{itemize}
        \item $\lim_{x \to -\infty} f'(x) = \lim_{x \to -\infty} (e^x + x) = -\infty$
        \item $\lim_{x \to +\infty} f'(x) = \lim_{x \to +\infty} (e^x + x) = +\infty$
    \end{itemize}
\end{rlist}

\Askhseis{\beta)} Η εξίσωση $e^x + x = \alpha$ (Ε) για τις διάφορες τιμές του πραγματικού αριθμού $\alpha$ γράφεται $f'(x) = \alpha$. Επειδή $f'(\mathbb{R}) = \mathbb{R}$ και η $f'$ είναι γνησίως αύξουσα έπεται ότι υπάρχει μοναδικό $\rho \in \mathbb{R}$ τέτοιο, ώστε $f'(\rho) = \alpha$. Επομένως για οποιαδήποτε τιμή του $\alpha \in \mathbb{R}$, η εξίσωση (Ε) έχει μοναδική λύση.

\Askhseis{\gamma)} Για κάθε $x \in \mathbb{R}$, έχουμε:
\[ g'(x) = (\alpha x - f(x) - \alpha \rho - f'(\rho))' = \alpha - f'(x) \]
\[ g'(x) > 0 \Leftrightarrow \alpha - f'(x) > 0 \Leftrightarrow \alpha > f'(x) \Leftrightarrow f'(\rho) > f'(x) \Leftrightarrow x < \rho \]
\[ g'(x) < 0 \Leftrightarrow x > \rho \]
\[ g'(\rho) = 0 \Leftrightarrow x = \rho \]

\begin{wrapfigure}[9]{r}{5.5cm}
    \begin{tabular}{|c|c|c|c|c|c|} \hline
        $x$ & \multicolumn{2}{c|}{$-\infty$} & $\rho$ & \multicolumn{2}{c|}{$+\infty$} \\ \hline
        $g'(x)$ & & {\color{maincolor} $+$} & $0$ & {\color{maincolor} $-$} & \\ \hline
        $g(x)$ & & {\color{maincolor} $\nearrow$} & OM & {\color{maincolor} $\searrow$} & \\ \hline
    \end{tabular}
\end{wrapfigure}
Η μονοτονία και τα ακρότατα της $g$ φαίνονται στον διπλανό πίνακα.

Οπότε η $g$ παρουσιάζει στη θέση $\rho$ μέγιστη τιμή, την:
\[ g(\rho) = \alpha \rho - f(\rho) \]

\Askhseis{β τρόπος}
Επειδή η συνάρτηση $f$ είναι κυρτή στο $\mathbb{R}$, τότε η εφαπτομένη της γραφικής παράστασης της $f$ στο σημείο $(\rho, f(\rho))$ βρίσκεται "κάτω" από τη γραφική παράσταση, με εξαίρεση το σημείο επαφής τους $(\rho, f(\rho))$.

Άρα για κάθε $x \in \mathbb{R}$ ισχύει:
\begin{align*}
    f(x) &\ge f'(\rho)(x - \rho) + f(\rho) \Leftrightarrow \\
    f(x) &\ge xf'(\rho) - \rho f'(\rho) + f(\rho) \Leftrightarrow \\
    \rho f'(\rho) - f(\rho) &\ge xf'(\rho) - f(x) \Leftrightarrow \\
    \alpha \rho - f(\rho) &\ge \alpha x - f(x) \Leftrightarrow \\
    g(x) &\le g(\rho) \quad \text{όπου } g(\rho) = \rho f'(\rho) - f(\rho)
\end{align*}
και έπεται το ζητούμενο.